\section{Homotopy and the shift functor}
\label{section-homotopy-shift}

\noindent
It is an annoying feature that signs and indices
have to be part of any discussion of homological
algebra\footnote{Please let us know if you notice sign errors or
if you have improvements to our conventions.}.

\begin{definition}
\label{definition-shift}
Let $\mathcal{A}$ be an additive category.
Let $A_\bullet$ be a chain complex
with boundary maps $d_{A, n} : A_n \to A_{n - 1}$.
For any $k \in \mathbf{Z}$ we define the
{\it $k$-shifted chain complex $A[k]_\bullet$}
as follows:
\begin{enumerate}
\item we set $A[k]_n = A_{n + k}$, and
\item we set $d_{A[k], n} : A[k]_n \to A[k]_{n - 1}$
equal to $d_{A[k], n} = (-1)^k d_{A, n + k}$.
\end{enumerate}
If $f : A_\bullet \to B_\bullet$ is a morphism of
chain complexes, then we let
$f[k] : A[k]_\bullet \to B[k]_\bullet$ be the
morphism of chain complexes with
$f[k]_n = f_{k + n}$.
\end{definition}

\noindent
Of course this means we have functors
$[k] : \text{Ch}(\mathcal{A}) \to \text{Ch}(\mathcal{A})$
which mutually commute (on the nose, without
any intervening isomorphisms of functors),
such that $A[k][l]_\bullet = A[k + l]_\bullet$ and
with $[0] = \text{id}_{\text{Ch}(\mathcal{A})}$.

\medskip\noindent
Recall that we view $\mathcal{A}$ as a full subcategory of
$\text{Ch}(\mathcal{A})$, see Section \ref{section-complexes}.
Thus for any object $A$ of $\mathcal{A}$ the notation $A[k]$
refers to the unique chain complex zero in all degrees except having
$A$ in degree $-k$.

\begin{definition}
\label{definition-homology-shift}
Let $\mathcal{A}$ be an abelian category.
Let $A_\bullet$ be a chain complex
with boundary maps $d_{A, n} : A_n \to A_{n - 1}$.
For any $k \in \mathbf{Z}$ we identify
{\it $H_{i + k}(A_\bullet) \rightarrow H_i(A[k]_\bullet)$}
via the identification
$A_{i + k} = A[k]_i$.
\end{definition}

\noindent
This identification is functorial in $A_\bullet$.
Note that since no signs are involved in this
definition we actually get a compatible system
of identifications of all the homology
objects $H_{i - k}(A[k]_\bullet)$, which are
further compatible with the identifications
$A[k][l]_\bullet = A[k + l]_\bullet$ and
with $[0] = \text{id}_{\text{Ch}(\mathcal{A})}$.

\medskip\noindent
Let $\mathcal{A}$ be an additive category.
Suppose that $A_\bullet$ and $B_\bullet$ are
chain complexes, $a, b : A_\bullet \to B_\bullet$ are
morphisms of chain complexes, and $\{h_i : A_i \to B_{i + 1}\}$
is a homotopy between $a$ and $b$. Recall that this means
that
$a_i - b_i = d_{i + 1} \circ h_i + h_{i - 1} \circ d_i$.
What if $a = b$? Then we obtain the formula
$0 = d_{i + 1} \circ h_i + h_{i - 1} \circ d_i$,
in other words, $ - d_{i + 1} \circ h_i = h_{i - 1} \circ d_i$.
By definition above this means the collection $\{h_i\}$
above defines a morphism of chain complexes
$$
A_\bullet \longrightarrow B[1]_\bullet.
$$
Such a thing is the same as a morphism $A[-1]_\bullet \to B_\bullet$
by our remarks above. This proves the following lemma.

\begin{lemma}
\label{lemma-homotopy-shift}
Let $\mathcal{A}$ be an additive category.
Suppose that $A_\bullet$ and $B_\bullet$ are
chain complexes. Given any morphism of chain
complexes $a : A_\bullet \to B_\bullet$ there
is a bijection between the set of homotopies
from $a$ to $a$ and
$\Mor_{\text{Ch}(\mathcal{A})}(A_\bullet, B[1]_\bullet)$.
More generally, the set of homotopies between
$a$ and $b$ is either empty or a principal homogeneous
space under the group
$\Mor_{\text{Ch}(\mathcal{A})}(A_\bullet, B[1]_\bullet)$.
\end{lemma}

\begin{proof}
See above.
\end{proof}

\begin{lemma}
\label{lemma-ses-termwise-split}
Let $\mathcal{A}$ be an abelian category.
Let
$$
0 \to A_\bullet \to B_\bullet \to C_\bullet \to 0
$$
be a short exact sequence of complexes.
Suppose that $\{s_n : C_n \to B_n\}$ is a family
of morphisms which split the short exact sequences
$0 \to A_n \to B_n \to C_n \to 0$. Let
$\pi_n : B_n \to A_n$ be the associated
projections, see Lemma \ref{lemma-ses-split}.
Then the family of morphisms
$$
\pi_{n - 1} \circ d_{B, n} \circ s_n
:
C_n \to A_{n - 1}
$$
define a morphism of complexes $\delta(s) : C_\bullet \to A[-1]_\bullet$.
\end{lemma}

\begin{proof}
Denote $i : A_\bullet \to B_\bullet$ and $q : B_\bullet \to C_\bullet$
the maps of complexes in the short exact sequence. Then
$i_{n - 1} \circ \pi_{n - 1} \circ d_{B, n} \circ s_n =
d_{B, n} \circ s_n - s_{n - 1} \circ d_{C, n}$. Hence
$i_{n - 2} \circ d_{A, n - 1} \circ \pi_{n - 1} \circ d_{B, n} \circ s_n =
d_{B, n - 1} \circ (d_{B, n} \circ s_n - s_{n - 1} \circ d_{C, n}) =
- d_{B, n - 1} \circ s_{n - 1} \circ d_{C, n}$ as desired.
\end{proof}

\begin{lemma}
\label{lemma-ses-termwise-split-long}
Notation and assumptions as in Lemma \ref{lemma-ses-termwise-split} above.
The morphism of complexes $\delta(s) : C_\bullet \to A[-1]_\bullet$
induces the maps
$$
H_i(\delta(s)) :
H_i(C_\bullet) \longrightarrow H_i(A[-1]_\bullet) = H_{i - 1}(A_\bullet)
$$
which occur in the long exact homology sequence associated
to the short exact sequence of chain complexes by
Lemma \ref{lemma-long-exact-sequence-chain}.
\end{lemma}

\begin{proof}
Omitted.
\end{proof}

\begin{lemma}
\label{lemma-ses-termwise-split-homotopy}
Notation and assumptions as in Lemma \ref{lemma-ses-termwise-split} above.
Suppose $\{s'_n : C_n \to B_n\}$ is a second choice of splittings.
Write $s'_n = s_n + i_n \circ h_n$ for some unique
morphisms $h_n : C_n \to A_n$. The family of maps
$\{h_n : C_n \to A[-1]_{n + 1}\}$ is a homotopy between
the associated morphisms
$\delta(s), \delta(s') : C_\bullet \to A[-1]_\bullet$.
\end{lemma}

\begin{proof}
Omitted.
\end{proof}

\begin{definition}
\label{definition-shift-cochain}
Let $\mathcal{A}$ be an additive category.
Let $A^\bullet$ be a cochain complex
with boundary maps $d_A^n : A^n \to A^{n + 1}$.
For any $k \in \mathbf{Z}$ we define the
{\it $k$-shifted cochain complex $A[k]^\bullet$}
as follows:
\begin{enumerate}
\item we set $A[k]^n = A^{n + k}$, and
\item we set $d_{A[k]}^n : A[k]^n \to A[k]^{n + 1}$
equal to $d_{A[k]}^n = (-1)^k d_A^{n + k}$.
\end{enumerate}
If $f : A^\bullet \to B^\bullet$ is a morphism of
cochain complexes, then we let
$f[k] : A[k]^\bullet \to B[k]^\bullet$ be the
morphism of cochain complexes with
$f[k]^n = f^{k + n}$.
\end{definition}

\noindent
Of course this means we have functors
$[k] : \text{CoCh}(\mathcal{A}) \to \text{CoCh}(\mathcal{A})$
which mutually commute (on the nose, without
any intervening isomorphisms of functors) and
such that $A[k][l]^\bullet = A[k + l]^\bullet$ and
with $[0] = \text{id}_{\text{CoCh}(\mathcal{A})}$.

\medskip\noindent
Recall that we view $\mathcal{A}$ as a full subcategory of
$\text{CoCh}(\mathcal{A})$, see Section \ref{section-complexes}.
Thus for any object $A$ of $\mathcal{A}$ the notation $A[k]$
refers to the unique cochain complex zero in all degrees except having
$A$ in degree $-k$.

\begin{definition}
\label{definition-cohomology-shift}
Let $\mathcal{A}$ be an abelian category.
Let $A^\bullet$ be a cochain complex
with boundary maps $d_A^n : A^n \to A^{n + 1}$.
For any $k \in \mathbf{Z}$ we identify
{\it $H^{i + k}(A^\bullet) \longrightarrow H^i(A[k]^\bullet)$}
via the identification $A^{i + k} = A[k]^i$.
\end{definition}

\noindent
This identification is functorial in $A^\bullet$.
Note that since no signs are involved in this
definition we actually get a compatible system
of identifications of all the homology
objects $H^{i - k}(A[k]^\bullet)$, which are
further compatible with the identifications
$A[k][l]^\bullet = A[k + l]^\bullet$ and
with $[0] = \text{id}_{\text{CoCh}(\mathcal{A})}$.

\medskip\noindent
Let $\mathcal{A}$ be an additive category.
Suppose that $A^\bullet$ and $B^\bullet$ are
cochain complexes, $a, b : A^\bullet \to B^\bullet$ are
morphisms of cochain complexes, and $\{h^i : A^i \to B^{i - 1}\}$
is a homotopy between $a$ and $b$. Recall that this means
that
$a^i - b^i = d^{i - 1} \circ h^i + h^{i + 1} \circ d^i$.
What if $a = b$? Then we obtain the formula
$0 = d^{i - 1} \circ h^i + h^{i + 1} \circ d^i$,
in other words, $ - d^{i - 1} \circ h^i = h^{i + 1} \circ d^i$.
By definition above this means the collection $\{h^i\}$
above defines a morphism of cochain complexes
$$
A^\bullet \longrightarrow B[-1]^\bullet.
$$
Such a thing is the same as a morphism $A[1]^\bullet \to B^\bullet$
by our remarks above. This proves the following lemma.

\begin{lemma}
\label{lemma-homotopy-shift-cochain}
Let $\mathcal{A}$ be an additive category.
Suppose that $A^\bullet$ and $B^\bullet$ are
cochain complexes. Given any morphism of cochain
complexes $a : A^\bullet \to B^\bullet$ there
is a bijection between the set of homotopies
from $a$ to $a$ and
$\Mor_{\text{CoCh}(\mathcal{A})}(A^\bullet, B[-1]^\bullet)$.
More generally, the set of homotopies between
$a$ and $b$ is either empty or a principal homogeneous
space under the group
$\Mor_{\text{CoCh}(\mathcal{A})}(A^\bullet, B[-1]^\bullet)$.
\end{lemma}

\begin{proof}
See above.
\end{proof}

\begin{lemma}
\label{lemma-ses-termwise-split-cochain}
Let $\mathcal{A}$ be an additive category.
Let
$$
0 \to A^\bullet \to B^\bullet \to C^\bullet \to 0
$$
be a complex (!) of complexes.
Suppose that we are given splittings $B^n = A^n \oplus C^n$
compatible with the maps in the displayed sequence.
Let $s^n : C^n \to B^n$ and $\pi^n : B^n \to A^n$ be the
corresponding maps. Then the family of morphisms
$$
\pi^{n + 1} \circ d_B^n \circ s^n
:
C^n \to A^{n + 1}
$$
define a morphism of complexes $\delta : C^\bullet \to A[1]^\bullet$.
\end{lemma}

\begin{proof}
Denote $i : A^\bullet \to B^\bullet$ and $q : B^\bullet \to C^\bullet$
the maps of complexes in the short exact sequence. Then
$i^{n + 1} \circ \pi^{n + 1} \circ d_B^n \circ s^n =
d_B^n \circ s^n - s^{n + 1} \circ d_C^n$. Hence
$i^{n + 2} \circ d_A^{n + 1} \circ \pi^{n + 1} \circ d_B^n \circ s^n =
d_B^{n + 1} \circ (d_B^n \circ s^n - s^{n + 1} \circ d_C^n) =
- d_B^{n + 1} \circ s^{n + 1} \circ d_C^n$ as desired.
\end{proof}

\begin{lemma}
\label{lemma-ses-termwise-split-long-cochain}
Notation and assumptions as in
Lemma \ref{lemma-ses-termwise-split-cochain} above.
Assume in addition that $\mathcal{A}$ is abelian.
The morphism of complexes $\delta : C^\bullet \to A[1]^\bullet$
induces the maps
$$
H^i(\delta) :
H^i(C^\bullet) \longrightarrow H^i(A[1]^\bullet) = H^{i + 1}(A^\bullet)
$$
which occur in the long exact homology sequence associated
to the short exact sequence of cochain complexes by
Lemma \ref{lemma-long-exact-sequence-cochain}.
\end{lemma}

\begin{proof}
Omitted.
\end{proof}

\begin{lemma}
\label{lemma-ses-termwise-split-homotopy-cochain}
Notation and assumptions as in
Lemma \ref{lemma-ses-termwise-split-cochain}.
Let $\alpha : A^\bullet \to B^\bullet$,
$\beta : B^\bullet \to C^\bullet$ be the given
morphisms of complexes.
Suppose $(s')^n : C^n \to B^n$ and $(\pi')^n : B^n \to A^n$
is a second choice of splittings.
Write $(s')^n = s^n + \alpha^n \circ h^n$ and
$(\pi')^n = \pi^n + g^n \circ \beta^n$ for some unique
morphisms $h^n : C^n \to A^n$ and $g^n : C^n \to A^n$. Then
\begin{enumerate}
\item $g^n = - h^n$, and
\item the family of maps $\{g^n : C^n \to A[1]^{n - 1}\}$ is a homotopy
between $\delta, \delta' : C^\bullet \to A[1]^\bullet$, more precisely
$(\delta')^n = \delta^n + g^{n + 1} \circ d_C^n + d_{A[1]}^{n - 1} \circ g^n$.
\end{enumerate}
\end{lemma}

\begin{proof}
As $(s')^n$ and $(\pi')^n$ are splittings we have $(\pi')^n \circ (s')^n = 0$.
Hence
$$
0 = ( \pi^n + g^n \circ \beta^n ) \circ ( s^n + \alpha^n \circ h^n ) =
g^n \circ \beta^n \circ s^n + \pi^n \circ \alpha^n \circ h^n =
g^n + h^n
$$
which proves (1). We compute $(\delta')^n$ as follows
$$
( \pi^{n + 1} + g^{n + 1} \circ \beta^{n + 1} )
\circ d_B^n \circ
( s^n + \alpha^n \circ h^n )
= \delta^n + g^{n + 1} \circ d_C^n + d_A^n \circ h^n
$$
Since $h^n = -g^n$ and since $d_{A[1]}^{n - 1} = -d_A^n$ we conclude that (2)
holds.
\end{proof}




